\paragraph{尾部可加闭包的 Ap\'ery 规约：二生成正规形、同余截断账本与复位时钟矩不变量}
本段在定理 \ref{thm:xi-window6-tail-semigroup-frobenius-threshold} 的尾部可加机制之上，把可达性从“生成元集合”进一步规约为一个同余阈值族，并给出与之等价的正规形、不可达集合的有限截断剖分、以及复位间隔时钟的矩不变量。
\par

\begin{theorem}[尾部可加机制的二生成规约与规范解]\label{thm:xi-window6-tail-semigroup-two-generator-normal-form}
在 window-\(6\) uplift 的尾部偏移可加闭包口径下，设可加生成元集合为 \(\{21,34,55\}\)。
由于 \(55=21+34\)，相应可达半群满足
\[
S_6=\langle 21,34,55\rangle=\langle 21,34\rangle\subset\mathbb N.
\]
对任意 \(n\in\mathbb N\)，令 \(r\equiv n\ (\mathrm{mod}\ 21)\)。
定义
\[
b_r:=13r\ (\mathrm{mod}\ 21)\in\{0,1,\dots,20\},
\qquad
w_r:=34\,b_r,
\]
则 \(n\in\langle 21,34\rangle\) 当且仅当 \(n\ge w_r\)。
并且在该可达情形下，存在唯一一组
\(
(a,b)\in\mathbb N\times\{0,1,\dots,20\}
\)
满足
\[
n=21a+34b,\qquad b=b_r,\qquad a=\frac{n-w_r}{21}.
\]
因此可达性在同余类 \(r\) 上具有常数时间判据，并诱导唯一的“模 \(21\) 规范形”（把 \(b\) 规约到 \(\{0,\dots,20\}\)）。
\end{theorem}
\leanverified{paper\_xi\_window6\_tail\_semigroup\_two\_generator\_normal\_form}

\begin{proof}
恒等式 \(55=21+34\) 说明 \(\langle 21,34,55\rangle=\langle 21,34\rangle\)。
又因 \(34\equiv 13\ (\mathrm{mod}\ 21)\) 且 \(13^2\equiv 1\ (\mathrm{mod}\ 21)\)，故 \(13\) 为 \(34\) 在模 \(21\) 下的逆元，从而
\(
34\,b_r\equiv r\ (\mathrm{mod}\ 21)
\)
成立，且 \(b_r\) 在 \(\{0,\dots,20\}\) 内唯一。
\par
若 \(n=21a+34b\) 可达，则必有 \(b\equiv b_r\ (\mathrm{mod}\ 21)\)，从而 \(b=b_r+21k\)（\(k\in\mathbb N\)）并推出
\(
n=w_r+21(a+34k)\ge w_r
\)。
反之，若 \(n\equiv r\ (\mathrm{mod}\ 21)\) 且 \(n\ge w_r\)，则 \(n-w_r\) 为 \(21\) 的非负倍数，取 \(a=(n-w_r)/21\ge 0\) 与 \(b=b_r\) 即得分解。
唯一性由 \(b\in\{0,\dots,20\}\) 的规约条件强制：若 \(21a+34b=21a'+34b'\) 且 \(b,b'\in\{0,\dots,20\}\)，则 \(34(b-b')\) 为 \(21\) 的倍数，因 \(\gcd(34,21)=1\) 得 \(b=b'\)，进而 \(a=a'\)。证毕。
\end{proof}

\begin{proposition}[Ap\'ery 阈值族与不可达计数的同余分解]\label{prop:xi-window6-tail-semigroup-apery-thresholds}
\leanverified{paper\_xi\_window6\_tail\_semigroup\_apery\_thresholds}
令 \(w_r\) 如定理 \ref{thm:xi-window6-tail-semigroup-two-generator-normal-form}。
则 \(\{w_r\}_{r=0}^{20}\) 构成 \(\langle 21,34\rangle\) 关于模 \(21\) 的 Ap\'ery 阈值族：
\[
w_r=\min\{s\in\langle 21,34\rangle:\ s\equiv r\ (\mathrm{mod}\ 21)\}.
\]
并且同余类 \(r\) 内的不可达点精确个数为
\[
g_r:=\frac{w_r-r}{21}\in\mathbb N,
\]
从而不可达集合可写为有限截断并
\(
\bigcup_{r=0}^{20}\{r,\,r+21,\dots,r+21(g_r-1)\}
\)。
其全量数据由表 \ref{tab:window6_tail_semigroup_apery_21_34} 给出，且总不可达计数满足
\[
\sum_{r=0}^{20}g_r=330.
\]
\end{proposition}

% AUTO-GENERATED by scripts/exp_window6_tail_semigroup_apery_21_34.py
\begin{table}[H]
\centering
\small
\setlength{\tabcolsep}{7pt}
\renewcommand{\arraystretch}{1.12}
\caption{Ap\'ery thresholds for the numerical semigroup $\langle 21,34\rangle$ modulo $21$: for each residue $r\in\{0,\dots,20\}$ we list the canonical $b_r\in\{0,\dots,20\}$ with $34b_r\equiv r\pmod{21}$, the Ap\'ery element $w_r=34b_r$, and the gap count $g_r=(w_r-r)/21$.}
\label{tab:window6_tail_semigroup_apery_21_34}
\begin{tabular}{r r r r}
\toprule
$r$ & $b_r$ & $w_r$ & $g_r$ \\
\midrule
0 & 0 & 0 & 0 \\
1 & 13 & 442 & 21 \\
2 & 5 & 170 & 8 \\
3 & 18 & 612 & 29 \\
4 & 10 & 340 & 16 \\
5 & 2 & 68 & 3 \\
6 & 15 & 510 & 24 \\
7 & 7 & 238 & 11 \\
8 & 20 & 680 & 32 \\
9 & 12 & 408 & 19 \\
10 & 4 & 136 & 6 \\
11 & 17 & 578 & 27 \\
12 & 9 & 306 & 14 \\
13 & 1 & 34 & 1 \\
14 & 14 & 476 & 22 \\
15 & 6 & 204 & 9 \\
16 & 19 & 646 & 30 \\
17 & 11 & 374 & 17 \\
18 & 3 & 102 & 4 \\
19 & 16 & 544 & 25 \\
20 & 8 & 272 & 12 \\
\bottomrule
\end{tabular}
\end{table}


\begin{proposition}[维数可达集合的有限截断同余并完全刻画]\label{prop:xi-window6-dimension-reachability-congruence-union}
令总维数写为 \(D=12+n\)（其中 \(n\) 为尾部可加增量）。
则在定理 \ref{thm:xi-window6-tail-semigroup-two-generator-normal-form} 的口径下：
\begin{enumerate}
  \item \(D\) 可达当且仅当存在 \(r\in\{0,\dots,20\}\) 使 \(n\equiv r\ (\mathrm{mod}\ 21)\) 且 \(n\in w_r+21\mathbb N\)；
  等价地，
  \[
  D\in\bigcup_{r=0}^{20}\bigl(12+w_r+21\mathbb N\bigr).
  \]
  \item \(D\) 不可达当且仅当存在 \(r\) 使 \(n=r+21k\) 且 \(0\le k<g_r\)；
  因此不可达集合是 \(21\) 条算术级数的有限截断并，其截断长度由 \(\{g_r\}\) 完全给出。
\end{enumerate}
\end{proposition}
\leanverified{paper\_xi\_window6\_dimension\_reachability\_congruence\_union}

\begin{proposition}[最小可达代表与 Zeckendorf 签名的单射编码]\label{prop:xi-window6-minrep-zeckendorf-signature-injection}
对每个 \(r\in\{0,\dots,20\}\)，定义同余类的最小可达维数
\[
D_r^{\min}:=12+w_r=12+34\,b_r.
\]
令 \(S(D)\subset\mathbb N\) 表示整数 \(D\) 的 Zeckendorf 指标集合（即 \(D=\sum_{k\in S(D)}F_k\)，且 \(S(D)\) 不含相邻指标）。
则 \(\{S(D_r^{\min})\}_{r=0}^{20}\) 两两不同；
并且对全部 \(r\)，有 \(2\in S(D_r^{\min})\) 且 \(3\notin S(D_r^{\min})\)。
其全量数据由表 \ref{tab:window6_tail_semigroup_minrep_zeckendorf_signatures} 给出。
因此，“同余类 \(r\)”在最小代表层面可被 \(S(D_r^{\min})\) 无歧义编码。
\end{proposition}
\leanverified{paper\_xi\_window6\_minrep\_zeckendorf\_signature\_injection}

% AUTO-GENERATED by scripts/exp_window6_tail_semigroup_apery_21_34.py
\begin{table}[H]
\centering
\small
\setlength{\tabcolsep}{7pt}
\renewcommand{\arraystretch}{1.12}
\caption{Minimal reachable representatives $D_r^{\min}=12+w_r$ for each residue class $r\pmod{21}$, together with their Zeckendorf index sets $S(D_r^{\min})$ in Fibonacci indices (so $D_r^{\min}=\sum_{k\in S}F_k$).}
\label{tab:window6_tail_semigroup_minrep_zeckendorf_signatures}
\begin{tabular}{r r l}
\toprule
$r$ & $D_r^{\min}$ & $S(D_r^{\min})$ \\
\midrule
0 & 12 & $\{2,4,6\}$ \\
1 & 454 & $\{2,8,10,14\}$ \\
2 & 182 & $\{2,4,9,12\}$ \\
3 & 624 & $\{2,7,15\}$ \\
4 & 352 & $\{2,6,8,11,13\}$ \\
5 & 80 & $\{2,4,8,10\}$ \\
6 & 522 & $\{2,12,14\}$ \\
7 & 250 & $\{2,4,7,13\}$ \\
8 & 692 & $\{2,5,8,10,15\}$ \\
9 & 420 & $\{2,6,9,14\}$ \\
10 & 148 & $\{2,4,12\}$ \\
11 & 590 & $\{2,7,10,12,14\}$ \\
12 & 318 & $\{2,6,8,10,13\}$ \\
13 & 46 & $\{2,4,6,9\}$ \\
14 & 488 & $\{2,8,11,14\}$ \\
15 & 216 & $\{2,4,7,10,12\}$ \\
16 & 658 & $\{2,7,9,15\}$ \\
17 & 386 & $\{2,6,14\}$ \\
18 & 114 & $\{2,4,8,11\}$ \\
19 & 556 & $\{2,9,12,14\}$ \\
20 & 284 & $\{2,4,7,9,13\}$ \\
\bottomrule
\end{tabular}
\end{table}


\begin{proposition}[同余置换的自对偶性：\texorpdfstring{$r\mapsto b_r$}{r->br} 为对合]\label{prop:xi-window6-congruence-involution}
映射 \(r\mapsto b_r=13r\ (\mathrm{mod}\ 21)\) 满足 \(13^2\equiv 1\ (\mathrm{mod}\ 21)\)，从而为自反置换：
\[
b_{b_r}=r.
\]
因此 Ap\'ery 阈值族满足互换律
\[
w_{b_r}=34\,b_{b_r}=34\,r.
\]
\end{proposition}
\leanverified{paper\_xi\_window6\_congruence\_involution}

\begin{proposition}[阈值族的一阶守恒：平均最小代表的闭式]\label{prop:xi-window6-apery-first-moment}
由于 \(r\mapsto b_r\) 为 \(\{0,\dots,20\}\) 上的置换，得
\[
\sum_{r=0}^{20}b_r=\sum_{k=0}^{20}k=210,
\qquad
\sum_{r=0}^{20}w_r=34\sum_{r=0}^{20}b_r=7140.
\]
从而
\[
\frac{1}{21}\sum_{r=0}^{20}w_r=340,
\qquad
\frac{1}{21}\sum_{r=0}^{20}D_r^{\min}=12+340=352.
\]
\end{proposition}
\leanverified{paper\_xi\_window6\_apery\_first\_moment}

\begin{proposition}[两间隔复位时钟的显式矩不变量]\label{prop:xi-window6-reset-gap-mgf-moments}
在结论 \ref{cor:xi-window6-reset-clock-gap-frequencies} 的设定下，
令复位间隔随机变量 \(G\) 取值于 \(\{34,55\}\)，且
\[
\mathbb P(G=55)=\frac{1}{\varphi},\qquad \mathbb P(G=34)=\frac{1}{\varphi^2}.
\]
则其矩母函数为
\[
M_G(t):=\mathbb E[e^{tG}]
=\frac{1}{\varphi}e^{55t}+\frac{1}{\varphi^2}e^{34t},
\]
并且均值与方差满足闭式
\[
\mathbb E[G]=\frac{55}{\varphi}+\frac{34}{\varphi^2}=13+21\varphi,
\qquad
\mathrm{Var}(G)=441(2\varphi-3).
\]
\end{proposition}

\begin{corollary}[Fibonacci 表达的漂移耦合]\label{cor:xi-window6-reset-gap-fib-coupling}
记 \(F_7=13\)、\(F_8=21\)。
则命题 \ref{prop:xi-window6-reset-gap-mgf-moments} 的漂移与涨落可重写为
\[
\mathbb E[G]=F_7+F_8\varphi,
\qquad
\mathrm{Var}(G)=F_8^2(2\varphi-3),
\]
从而复位漂移与涨落在整数层面由 \((F_7,F_8)\) 与同一无理常数 \(\varphi\) 完全确定。
\end{corollary}

\begin{proposition}[不可达集合的质量不变量与二阶一致性账本]\label{prop:xi-window6-gap-ledger-invariants}
令 \(\mathcal G\subset\mathbb N\) 为 \(\langle 21,34\rangle\) 的不可达增量集合（即命题 \ref{prop:xi-window6-tail-semigroup-apery-thresholds} 的有限截断并对应点集）。
则其基数、总权重以及由截断长度 \(\{g_r\}\) 导出的二阶组合量满足严格恒等式
% AUTO-GENERATED by scripts/exp_window6_tail_semigroup_gap_ledger_21_34.py
\begin{equation}\label{eq:window6_tail_semigroup_gap_ledger_21_34}
\begin{aligned}
|\mathcal{G}|&=330,
&\sum_{n\in\mathcal{G}} n&=75460,
&\frac{1}{|\mathcal{G}|}\sum_{n\in\mathcal{G}} n&=\frac{686}{3},\\
\sum_{r=0}^{20} g_r&=330,
&\sum_{r=0}^{20}\binom{g_r}{2}&=3432,
&\sum_{r=0}^{20} g_r^2&=7194,
&\sum_{r=0}^{20} r\,g_r&=3388.
\end{aligned}
\end{equation}

因此，维数层面的不可达点集 \(\{12+n:\ n\in\mathcal G\}\) 的均值为 \(12+\frac{686}{3}=\frac{722}{3}\)。
上述恒等式与表 \ref{tab:window6_tail_semigroup_apery_21_34} 的阈值族共同构成一组可核验的一致性账本；其任一破坏都等价于同余阈值族或“尾部可加”机制假设的失配。
\end{proposition}
\leanverified{paper\_xi\_window6\_gap\_ledger\_invariants}

\paragraph{Sturmian 复位时钟的有界漂移与常数时间可计算性}

\begin{proposition}[两间隔 Sturmian 复位时钟的有界漂移：抖动上界等于间隔差]\label{prop:xi-window6-reset-clock-bounded-drift}
\leanverified{paper\_xi\_window6\_reset\_clock\_bounded\_drift}
令 window-\(6\) 的复位事件序列如定理 \ref{thm:terminal-reset-events-sturmian}，并记相邻复位间隔为
\[
G_n:=r_{n+1}-r_n\in\{34,55\}\qquad(n\ge 1).
\]
定义部分和（复位时刻增量）
\[
T_n:=\sum_{j=1}^{n}G_j\qquad(n\ge 1),
\]
并令均值漂移常数
\[
\mu:=\mathbb E[G]=\frac{55}{\varphi}+\frac{34}{\varphi^2}=13+21\varphi
\]
如命题 \ref{prop:xi-window6-reset-gap-mgf-moments}。
则存在截距参数 \(\theta\in[0,1)\)，使得长间隔计数
\[
S_n:=\#\{1\le j\le n:\ G_j=55\}
\]
满足 \(S_n=\lfloor n/\varphi+\theta\rfloor\)，并且漂移偏差具有统一有界性
\[
T_n-n\mu=21\Bigl(S_n-\frac{n}{\varphi}\Bigr),
\qquad
\bigl|T_n-n\mu\bigr|\le 21
\quad(\forall n\ge 1).
\]
因此该复位时钟在所有尺度上具有确定性的“零扩散”抖动：偏离均值线性漂移的误差不随 \(n\) 增长，且其最坏幅度恰等于两间隔之差 \(55-34=21\)。
\end{proposition}

\begin{proof}
由定理 \ref{thm:terminal-reset-events-sturmian}，指示序列 \(\ind_{\{G_j=55\}}\) 为 slope \(\varphi^{-1}\) 的 Sturmian 序列（等价于无理旋转的区间窗读出）。
Sturmian 序列的平衡性给出：存在 \(\theta\in[0,1)\) 使得
\(
S_n=\lfloor n/\varphi+\theta\rfloor
\)
且 \(|S_n-n/\varphi|\le 1\)。
另一方面，恒等式 \(T_n=55S_n+34(n-S_n)=34n+21S_n\) 与 \(\mu=34+21/\varphi\) 推出
\(
T_n-n\mu=21(S_n-n/\varphi)
\)；
由 \(|S_n-n/\varphi|\le 1\) 得 \(|T_n-n\mu|\le 21\)。证毕。
\end{proof}

\begin{corollary}[复位时钟的常数时间可计算性：Beatty 取整与无扫描更新]\label{cor:xi-window6-reset-clock-O1-computable}
在命题 \ref{prop:xi-window6-reset-clock-bounded-drift} 的记号下，给定 \(n\) 可在常数时间计算 \(T_n\)：
令 \(S_n=\lfloor n/\varphi+\theta\rfloor\)，则
\[
T_n=34n+21S_n.
\]
从而复位时刻集合 \(\{T_n\}\) 可由一次 Beatty 取整实现，不需枚举 \((G_1,\dots,G_n)\)。
\end{corollary}

\begin{remark}[漂移余差的旋转表示与同余不变量]\label{rem:xi-window6-reset-clock-residual-rotation-congruence}
由 \(S_n=\lfloor n/\varphi+\theta\rfloor\) 得
\[
S_n-\frac{n}{\varphi}=\theta-\Bigl\{\frac{n}{\varphi}+\theta\Bigr\}\in(\theta-1,\theta],
\qquad
T_n-n\mu=21\Bigl(\theta-\Bigl\{\frac{n}{\varphi}+\theta\Bigr\}\Bigr),
\]
其中 \(\{\cdot\}\) 为小数部分。
由于 \(1/\varphi\) 无理，\(\{\frac{n}{\varphi}+\theta\}\) 在 \([0,1)\) 上稠密，从而漂移余差 \(\{T_n-n\mu\}_{n\ge 1}\) 在某个长度为 \(21\) 的区间上稠密。
另一方面，\(\theta-\{\frac{n}{\varphi}+\theta\}\) 的符号仅由
\(\{\frac{n}{\varphi}+\theta\}\in[0,\theta]\) 或 \((\theta,1)\) 两端区间决定；
因此漂移余差的\emph{符号语义}在二元预算下即可压缩（对应 Sturmian 二元窗读出下的局部模式二分）。
另一方面，由 \(34\equiv 55\equiv 13\ (\mathrm{mod}\ 21)\) 得
\[
T_n\equiv 13n\ (\mathrm{mod}\ 21),
\]
故复位时刻在模 \(21\) 的同余类上呈现完全确定的有限剩余类结构。
\end{remark}

\begin{corollary}[复位时钟与尾部半群的可达性耦合：有限障碍仅发生于浅层]\label{cor:xi-window6-reset-clock-semigroup-coupling}
令 \(S=\langle 21,34\rangle\)。
则对任意整数 \(t\ge 715\)，存在 \(n\in\mathbb N\) 与 \(s\in S\) 使
\[
t=T_n+s
\qquad\text{且}\qquad
\bigl|t-(n\mu+s)\bigr|\le 21,
\]
其中 \(\mu\) 与 \(\{T_n\}\) 如命题 \ref{prop:xi-window6-reset-clock-bounded-drift}。
因此在深尺度上，“复位时钟（Sturmian）”与“尾部可达半群（Ap\'ery）”形成确定性的近共终结构；真正的不可达障碍仅可能出现在有限多个浅层尺度上。
\end{corollary}

\begin{proof}
由定理 \ref{thm:xi-window6-tail-semigroup-two-generator-normal-form}（或等价地由 \(\gcd(21,34)=1\) 的 Frobenius 余有限性），有 \(n\ge 660\Rightarrow n\in S\)。
取任意 \(t\ge 715\)。
注意 \(T_1\in\{34,55\}\)，故 \(s:=t-T_1\ge 660\)，从而 \(s\in S\) 且 \(t=T_1+s\)。
另一方面由命题 \ref{prop:xi-window6-reset-clock-bounded-drift}（取 \(n=1\)）得
\(
|T_1-\mu|\le 21
\)，于是
\(
|t-(\mu+s)|=|T_1-\mu|\le 21
\)。
证毕。
\end{proof}

\paragraph{尾部三元组与审计素数的算术闭环}

\begin{corollary}[GUT 三元尾部偏移与最大退化三元组的双源同一：\texorpdfstring{$\{21,34,55\}$}{\{21,34,55\}} 的自洽闭环]\label{cor:xi-tail-offsets-maxfiber-triple-closure}
在 window-\(6\) uplift 的尾部偏移口径下，非零偏移仅可能取自
\[
\Delta_{\mathrm{tail}}(6)\subseteq \{F_8,F_9,F_{10}\}=\{21,34,55\}
\]
（定理 \ref{thm:xi-window6-tail-semigroup-frobenius-threshold}）。
另一方面，由最大纤维闭式（定理 \ref{thm:pom-max-fiber}）在偶数分辨率上有 \(D_{2k}=F_{k+2}\)，从而
\[
D_{12}=F_8=21,\qquad D_{14}=F_9=34,\qquad D_{16}=F_{10}=55.
\]
因此同一 Fibonacci 三元组 \(\{21,34,55\}\) 以两条机制同时出现：一条来自 window-\(6\) 的尾部跃迁字母表，另一条来自更高分辨率的最大退化规模；该双源同一在纯离散层面封闭了“统一三分支”的自洽核对链条。
\end{corollary}

\begin{proof}
第一部分由定理 \ref{thm:xi-window6-tail-semigroup-frobenius-threshold}。
第二部分由定理 \ref{thm:pom-max-fiber} 取 \(k=6,7,8\) 并代入 \(F_8=21,F_9=34,F_{10}=55\)。
合并即得。证毕。
\end{proof}

\begin{corollary}[最大纤维的奇偶闭式与最大达到者 hidden-bit 分裂的有限态谱]\label{cor:xi-max-fiber-parity-hiddenbit-finite-spectrum}
令 \(D_m:=\max_{x\in X_m}d_m(x)\)。
则最大纤维满足奇偶闭式
\[
D_{2k}=F_{k+2},\qquad D_{2k+1}=2F_{k+1}
\qquad(k\ge 1)
\]
（定理 \ref{thm:pom-max-fiber}）。
对任意最大达到者 \(x\in X_m\)（即 \(d_m(x)=D_m\)），按 hidden-bit 分裂定义 \((d_{m,0}(x),d_{m,1}(x))\)（注 \ref{rem:pom-max-fiber-achievers-bsplit}），则在稳定阈值区间内分裂对满足闭式有限态结构：
\begin{itemize}
  \item \textbf{偶数 \(m=2k\ge 10\)：}
  \[
  (d_{2k,0}(x),d_{2k,1}(x))\in\{(F_{k+1},F_k),(F_k,F_{k+1})\}.
  \]
  \item \textbf{奇数 \(m=2k+1\ge 13\)：}
  \[
  (d_{2k+1,0}(x),d_{2k+1,1}(x))\in\{(F_{k+1},F_{k+1}),\ (F_{k+1}\pm F_{k-2},\,F_{k+1}\mp F_{k-2})\}.
  \]
\end{itemize}
特别地，对奇数窗 \(m=2k+1\ge 13\) 定义
\(
p_1(x):=d_{m,1}(x)/D_m
\)。
则三态谱满足
\[
p_1(x)\in\left\{\frac12,\ \frac12\pm\frac{F_{k-2}}{2F_{k+1}}\right\},
\qquad
\Delta_k:=\frac{F_{k-2}}{2F_{k+1}}=\frac12\varphi^{-3}+O(\varphi^{-2k}),
\]
从而 \(p_1\) 的极限谱为 \(\{\varphi^{-2},\tfrac12,\varphi^{-1}\}\)，并以指数速度收敛。
\end{corollary}

\begin{proof}
最大纤维的奇偶闭式由定理 \ref{thm:pom-max-fiber} 给出。
达到者分裂对的闭式有限态结构见注 \ref{rem:pom-max-fiber-achievers-bsplit}，而三态谱与极限谱见推论 \ref{cor:pom-max-fiber-achievers-hidden-bit-prob-spectrum}。
最后的渐近由 Binet 公式给出：\(F_{k-2}/F_{k+1}=\varphi^{-3}+O(\varphi^{-2k})\)。证毕。
\end{proof}

\begin{corollary}[审计素数 \texorpdfstring{$p_\star=571$}{p*=571} 的尾部半群 Ap\'ery 正规形与唯一性]\label{cor:xi-audit-prime-571-apery-normal-form}
令 \(p_\star:=571\) 为四块商图的树数指纹素数（推论 \ref{cor:xi-window6-block-tree-number-audit-prime}），并令尾部半群 \(S=\langle 21,34\rangle\)。
则 \(p_\star\in S\)，且在定理 \ref{thm:xi-window6-tail-semigroup-two-generator-normal-form} 的模 \(21\) 规范形约束 \(b\in\{0,\dots,20\}\) 下具有唯一表示
\[
p_\star=21\cdot 11+34\cdot 10.
\]
其中系数 \(b=10\) 正是同余类 \(p_\star\ (\mathrm{mod}\ 21)\) 的 Ap\'ery 正规形系数（即 \(34\cdot 10\) 为该同余类在 \(S\) 中的最小代表）。
\end{corollary}

\begin{proof}
计算 \(p_\star\equiv 571\equiv 4\ (\mathrm{mod}\ 21)\)。
由定理 \ref{thm:xi-window6-tail-semigroup-two-generator-normal-form}，规范系数为
\(
b_{4}=13\cdot 4\ (\mathrm{mod}\ 21)=10
\)，从而同余类阈值 \(w_4=34\,b_4=340\)。
于是
\(
p_\star-w_4=571-340=231=21\cdot 11
\)，得表示 \(571=21\cdot 11+34\cdot 10\)。
唯一性由同一定理的“模 \(21\) 规范形”唯一性直接推出。证毕。
\end{proof}
